% !PRINTED 93
\weberpage{93}
%<<CONT

\[
A=\begin{vmatrix}
 a_1^{(1)}&a_2^{(1)}&\cdots&a_n^{(1)}\\
 a_1^{(2)}&a_2^{(2)}&\cdots&a_n^{(2)}\\
 \cdots&\cdots&&\cdots\\
 a_1^{(n)}&a_2^{(n)}&\cdots&a_n^{(n)}
\end{vmatrix},\qquad
B=\begin{vmatrix}
 b_1^{(1)}&b_2^{(1)}&\cdots&b_n^{(1)}\\
 b_1^{(2)}&b_2^{(2)}&\cdots&b_n^{(2)}\\
 \cdots&\cdots&&\cdots\\
 b_1^{(n)}&b_2^{(n)}&\cdots&b_n^{(n)}
\end{vmatrix},
\]
and if $A_i^{(k)},B_i^{(k)}$ denote the first subdeterminants, one obtains
\begin{align}
 A x_k&=\sum_i A_k^{(i)}y_i,\tag{3}\\
 B y_i&=\sum_h B_h^{(i)}z_h,\tag{4}
\end{align}
and, by substitution of (4) in (3),
\begin{equation}
 ABx_k=\sum_h z_h\sum_i A_k^{(i)}B_h^{(i)}.\tag{5}
\end{equation}
But one may also proceed by inserting from (1) the expressions for $y$ into (2), whereby, if
\begin{equation}
 c_i^{(h)}=\sum_s a_i^{(s)}b_h^{(s)}\tag{6}
\end{equation}
is put, one obtains
\begin{equation}
 \sum_i c_i^{(h)}x_i=z_h.\tag{7}
\end{equation}
Now put
\[
C=\begin{vmatrix}
 c_1^{(1)}&c_2^{(1)}&\cdots&c_n^{(1)}\\
 c_1^{(2)}&c_2^{(2)}&\cdots&c_n^{(2)}\\
 \cdots&\cdots&&\cdots\\
 c_1^{(n)}&c_2^{(n)}&\cdots&c_n^{(n)}
\end{vmatrix}\tag{8}
\]
If $C_i^{(k)}$ denotes the subdeterminants of $C$, the solution of (7) gives
\begin{equation}
 Cx_k=\sum_h z_h C_k^{(h)},\tag{9}
\end{equation}
and comparison of (5) with (9) gives
\begin{equation}
 AB=C,\tag{10}
\end{equation}
and for the subdeterminants
\begin{equation}
 C_k^{(h)}=\sum_i A_k^{(i)}B_h^{(i)}.\tag{11}
\end{equation}
In this conclusion, however, there is still a gap, caused by the doubt whether (5) and (9) might not differ by a factor common to both sides. To remove this doubt we shall prove equation (10) directly; the correctness of (11) then follows.
